\subsection{Confidence sequence cho quantile}

Trong thí nghiệm này, ta khảo sát confidence sequence cho quantile
trên phân phối Cauchy chuẩn. Phân phối Cauchy được chọn vì đây là
một phân phối có đuôi nặng, nên các quan sát cực trị có thể xuất
hiện với giá trị rất lớn. Tuy nhiên, các quantile của phân phối
Cauchy vẫn tồn tại, do đó phương pháp dựa trên quantile vẫn có thể
được áp dụng.

Ta sinh tuần tự $T=100000$ quan sát độc lập từ phân phối
Cauchy chuẩn. Tại mỗi thời điểm $t$, ta chỉ sử dụng $t$ quan sát
đầu tiên để xây dựng khoảng tin cậy. Trong thí nghiệm này, ta chọn
mức ý nghĩa $\alpha=0.05$, tương ứng với mức tin cậy
$1-\alpha=0.95$.

Trước hết, ta xét bài toán xây dựng confidence sequence cho một
quantile cố định, cụ thể là quantile mức $p=0.9$. Với trường hợp này,
độ rộng sai số được chọn là
\begin{equation}
u_t=\sqrt{\frac{0.73\log\log(2.04t)+0.52\log(9.97/\alpha)}{t}}.
\end{equation}
Khi đó, tại mỗi thời điểm $t$, khoảng confidence sequence cho
quantile $Q(0.9)$ được cho bởi
\begin{equation}
\left[\widehat Q_t(0.9-u_t),\widehat Q_t(0.9+u_t)\right].
\end{equation}

Bên cạnh confidence sequence cho một quantile cố định, ta cũng xét
trường hợp mạnh hơn, trong đó khoảng tin cậy cần đúng đồng thời cho
mọi thời điểm $t$ và mọi mức quantile $p\in(0,1)$. Khi đó, độ rộng
sai số được chọn là
\begin{equation}
u_t=\sqrt{\frac{\log\log(et)+0.75\log(34/\alpha)}{t}}.
\end{equation}

\figurename~\ref{fig:confidence-sequence} minh họa hai kết quả chính của thí
nghiệm. Ở phía bên trái, hình biểu diễn hai cận của confidence
sequence cho quantile $0.9$ theo số mẫu $t$. Đường chấm ngang biểu
diễn quantile thật $Q(0.9)$. Ta thấy hai cận dần thu hẹp quanh
quantile thật khi số mẫu tăng. Ở phía bên phải, hình biểu diễn
confidence band cho toàn bộ hàm phân phối tích lũy tại ba thời điểm
$t=100$, $t=1000$ và $t=10000$. Khi $t$ tăng, confidence band hẹp
dần và tiến gần hơn đến CDF thật.

\begin{figure}[htbp]
    \centering
    \includegraphics[width=0.95\textwidth]{confidence-sequence.png}
    \caption{Confidence sequence cho quantile $0.9$ và confidence
    band cho toàn bộ CDF trên dữ liệu sinh từ Cauchy chuẩn.}
    \label{fig:confidence-sequence}
\end{figure}

\tablename~\ref{tab:quantile-coverage} trình bày kết quả định lượng cho
quantile $0.9$ tại các mốc $t=100$, $t=1000$, $t=10000$ và
$t=100000$. Cột cuối cho biết khoảng confidence sequence tại thời
điểm tương ứng có chứa quantile thật hay không. Tại tất cả các mốc
được kiểm tra, quantile thật $Q(0.9)$ đều nằm trong khoảng
confidence sequence. Khi số mẫu tăng từ $100$ lên $100000$, độ rộng
của khoảng giảm rõ rệt. Ở $t=100$, khoảng còn khá rộng, từ $1.0766$
đến $10.6478$. Đến $t=100000$, khoảng đã thu hẹp còn từ $2.8514$
đến $3.2892$.

\begin{table}[htbp]
  \centering
  \caption{Kiểm tra định lượng confidence sequence cho quantile
  $0.9$ của phân phối Cauchy$(0,1)$.}
  \label{tab:quantile-coverage}

  \begin{tblr}{
    width=\textwidth,
    colspec={cccccc},
    rowhead = 1,
    hlines,
    vlines
  }
    $t$ & $u_t$ & cận dưới & true $Q(0.9)$ & cận trên & covered \\
    100 & 0.1993 & 1.0766 & 3.0777 & 10.6478 & True \\
    1000 & 0.0651 & 1.8955 & 3.0777 & 11.0662 & True \\
    10000 & 0.0210 & 2.5256 & 3.0777 & 3.9465 & True \\
    100000 & 0.0068 & 2.8514 & 3.0777 & 3.2892 & True \\
  \end{tblr}
\end{table}

Ngoài việc kiểm tra trên một chuỗi dữ liệu cố định, ta cũng lặp lại
thí nghiệm nhiều lần để ước lượng xác suất bao phủ. Với mỗi lần lặp,
ta sinh một chuỗi dữ liệu mới và kiểm tra xem khoảng confidence
sequence có chứa quantile thật tại tất cả các thời điểm trong lưới
$t$ hay không. Tỉ lệ các lần kiểm tra thành công được dùng làm ước
lượng thực nghiệm cho xác suất bao phủ.

Đối với trường hợp một quantile cố định, ta kiểm tra quantile mức
$p=0.9$ với $100$ mốc số mẫu t. Đối với trường hợp đồng thời cho mọi
quantile, ta không thể kiểm tra trực tiếp vô hạn mọi mức $p\in(0,1)$ và
mọi thời điểm $t\in\mathbb{N}$. Do đó, trong mô phỏng, ta kiểm tra trên
một lưới hữu hạn gồm $99$ mức quantile từ $0.01$ đến $0.99$ và $50$ mốc
số mẫu $t$.

\tablename~\ref{tab:coverage-estimation} trình bày kết quả ước lượng bao phủ.
Cả hai trường hợp đều đạt bao phủ thực nghiệm bằng $1.0$, tức là cả $300$ lần
lặp đều thành công. Kết quả này lớn hơn mức bao phủ mục tiêu $0.95$.

\begin{table}[htbp]
  \centering
  \caption{Ước lượng xác suất bao phủ bằng $300$ lần mô phỏng
  độc lập trên phân phối Cauchy$(0,1)$.}
  \label{tab:coverage-estimation}

  \begin{tblr}{
    width=\textwidth,
    colspec={X[1.6,l] X[1,c] X[1,c] X[1,c]},
    rowhead = 1,
    hlines,
    vlines
  }
    Trường hợp & target & estimated & success / trials \\
    Fixed quantile $p=0.9$ & 0.95 & 1.00 & 300 / 300 \\
    All quantiles & 0.95 & 1.00 & 300 / 300 \\
  \end{tblr}
\end{table}

Kết quả trên phù hợp với bảo đảm lý thuyết của confidence sequence.
Tại các mốc số mẫu được kiểm tra, quantile thật $Q(0.9)$ đều nằm
trong khoảng tin cậy. Khi số mẫu $t$ tăng thì $u_t$ giảm, nên khoảng
tin cậy hẹp dần. Ban đầu cận trên có thể lớn do phân phối Cauchy có
đuôi nặng, nhưng khi có thêm dữ liệu, khoảng tin cậy tiến gần hơn
quanh quantile thật.

Kết quả mô phỏng lặp lại cũng cho thấy xác suất bao phủ thực nghiệm
không thấp hơn mức mục tiêu $1-\alpha=0.95$. Cụ thể, cả trường hợp
fixed quantile và all quantiles đều đạt $300/300$ lần thành công
trong thí nghiệm này.
